\section{Zorn's lemma}

\import{order/hausdorff.tex}
\import{order/bounds.tex}

The general form is stated for an arbitrary partial order.  The
subset-family form below is a corollary for developments where upper
bounds are supplied by unions.

\begin{theorem}\label{zorn_partial_order}
    Suppose $R$ is an order on $P$.
    Suppose for every $C$ such that $C$ is a chain in $P$ with respect to $R$
        there exists $u$ such that $u$ is an upper bound of $C$ in $P$
            with respect to $R$.
    Then there exists $m$ such that $m$ is a maximal element of $P$
        with respect to $R$.
\end{theorem}
\begin{proof}
    Let $Q=\tostrictorder{R}$.
    $R$ is an order by \obv{order_on_implies_order}.
    Thus $Q$ is a strict order by \obv{strictorder_from_order}.
    $R\subseteq P\times P$.
    $Q\subseteq R$ by \auto{tostrictorder,subseteq}.
    Hence $Q\subseteq P\times P$ by \obv{subseteq_transitive}.
    Take $C$ such that $C$ is a maximal chain in $P$ with respect to $Q$
        by \obv{hausdorff_maximal_chain}.
    Then $C$ is a maximal chain in $P$ with respect to $R$
        by \obv{maximal_chain_tostrictorder_iff}.
    Thus $C$ is a chain in $P$ with respect to $R$
        by \obv{maximal_chain_in_relation}.
    Hence $C\subseteq P$ and for all $a,b\in C$ we have
        $a=b$ or $a\mathrel{R}b$ or $b\mathrel{R}a$
        by \obv{chain_in_relation}.
    Take $u$ such that $u$ is an upper bound of $C$ in $P$
        with respect to $R$.

    Show that $u$ is a maximal element of $P$ with respect to $R$.
    \begin{subproof}
        $u\in P$ by \obv{upper_bound}.
        Show that for all $x\in P$ such that $u\mathrel{R}x$ we have $x=u$.
        \begin{subproof}
            Fix $x$.
            Assume $x\in P$ and $u\mathrel{R}x$.
            Show that $\cons{x}{C}$ is a chain in $P$ with respect to $R$.
            \begin{subproof}
                $C\subseteq P$.
                Hence $\cons{x}{C}\subseteq P$ by \obv{cons_subseteq_intro}.
                Show that for all $a,b\in\cons{x}{C}$ we have
                    $a=b$ or $a\mathrel{R}b$ or $b\mathrel{R}a$.
                \begin{subproof}
                    Fix $a,b\in\cons{x}{C}$.
                    \begin{byCase}
                    \caseOf{$a=x$.}
                        Follows by
                            \auto{cons,upper_bound,transitive}.
                    \caseOf{$b=x$.}
                        Follows by
                            \auto{cons,upper_bound,transitive}.
                    \caseOf{$a,b\in C$.}
                        Follows.
                    \end{byCase}
                \end{subproof}
                Follows by \obv{chain_in_relation}.
            \end{subproof}
            Thus $\cons{x}{C}$ is a chain in $P$ with respect to $Q$
                by \obv{chain_tostrictorder_iff}.
            $C\subseteq\cons{x}{C}$ by \obv{subseteq_cons_self}.
            Hence $\cons{x}{C}=C$ by
                \obv{maximal_chain_in_relation}.
            Thus $x\in C$ by \obv{cons_left}.
            Hence $x=u$ or $x\mathrel{R}u$ by \obv{upper_bound_elim}.
            Follows by \obv{antisymmetric}.
        \end{subproof}
        Follows by \obv{maximal_element}.
    \end{subproof}
\end{proof}

\begin{theorem}\label{zorn_subset_family}
    Suppose $P$ is inhabited.
    Suppose for every $C$ such that $C\subseteq P$ and $C$ is inhabited
        and $C$ is an inclusion chain we have $\unions{C}\in P$.
    Then there exists $M\in P$ such that for every $X\in P$ such that
        $M\subseteq X$ we have $X=M$.
\end{theorem}
\begin{proof}
    Let $R=\subseteqrel{P}$.
    $R$ is an order on $P$ by \auto{subseteqrel_is_order}.
    Show that for every $C$ such that $C$ is a chain in $P$ with respect to $R$
        there exists $u$ such that $u$ is an upper bound of $C$ in $P$
            with respect to $R$.
    \begin{subproof}
        Fix $C$.
        Assume $C$ is a chain in $P$ with respect to $R$.
        $C\subseteq P$ and $C$ is an inclusion chain
            by \auto{inclusion_chain_in_family_iff}.
        \begin{byCase}
        \caseOf{$C$ is inhabited.}
            $\unions{C}\in P$.
            For all $X\in C$ we have
                $X=\unions{C}$ or $X\mathrel{R}\unions{C}$
                by \auto{unions_iff,subseteq,subseteqrel,pair_eq_iff}.
            Thus $\unions{C}$ is an upper bound of $C$ in $P$
                with respect to $R$ by \obv{upper_bound}.
        \caseOf{$C$ is not inhabited.}
            Take $u$ such that $u\in P$.
            $u$ is an upper bound of $C$ in $P$ with respect to $R$
                by \auto{upper_bound}.
        \end{byCase}
    \end{subproof}
    Take $M$ such that $M$ is a maximal element of $P$ with respect to $R$
        by \obv{zorn_partial_order}.
    $M\in P$ by \obv{maximal_element}.
    For every $X\in P$ such that $M\subseteq X$ we have $X=M$
        by \obv{maximal_element,subseteqrel,pair_eq_iff}.
\end{proof}
